% Section B.5: the answers to the parts of lectures/L05/appendix/b_exercises.md that are not code.
% Every testbench message quoted was produced by a design built to Chapters 1 to 5 and broken in
% exactly the way the part describes.

\section{Chapter 5: The Register Bank}
\label{sol:c5}
The a) parts of both exercises are code: \code{uart_regs}, checked by \code{uart_regs_tb}, and the
completed \code{uart_top}, checked by \code{uart_top_tb}.

\solution{c5:ex:1}{\code{uart_regs}}
\ans{b} Advancing the FIFO ``whenever that index is read'' means decoding a pop from the address
alone, since the bus has no read strobe: \code{rx_pop} is high whenever \code{reg_idx} is
\code{REG_RX_DATA}. The run then fails one check earlier than you might expect, in Case~5 rather
than Case~6:

\begin{console}
uart_regs_tb: RX_DATA mismatch, got 0x00!
\end{console}

\noindent \code{bus_set_addr} puts the address on the bus, waits for a clock edge and then samples.
That edge pops the FIFO, so by the time the bench reads \code{RX_DATA} the byte it pushed is gone
and the front shows an entry that was never written. The ``a bare \code{RX_DATA} read must not pop''
check in Case~6 would fail too, but the run never gets there.

The system testbench, by contrast, \emph{passes} with the change. The bridge latches
\code{reg_rdata} on the same edge as the first pop, so it captures the byte before it disappears,
and \code{uart_top_tb} reads \code{RX_DATA} once and never looks again. What it does not see is the
rest: the bridge holds \code{reg_addr} at the last command's index until the next command, so the
pop repeats on every clock and an eight-entry FIFO is empty 160\,ns after the command byte.

That is the general problem, and the abort rule makes it plain. A read with a side effect has to
decide \emph{when} the side effect happens. If the pop happens when the value is latched, after the
command byte, then an aborted read, \code{SS} rising before the fifth byte, has discarded a byte the
master never received, and the protocol promises that an abort has no side effects. If the pop waits
for the fifth byte, it is a commit on completion, the discipline a write has; but the bridge commits
only writes, with \code{reg_write}, and has no ``read completed'' strobe to give the bank, so both
the bridge and the bank would have to change. The pure read plus a separate \code{RX_POP} needs
none of that: reading has no side effect to abort, and the pop is a write, which the bridge already
commits exactly once, on completion, and never on an abort.

\ans{c} \code{STATUS} bit 0, TX ready, is \code{not tx_full_s}, from the TX FIFO. Bit 1, RX valid, is
\code{not rx_empty_s}, from the RX FIFO. Bit 2, Error, is the OR of the three \code{err_flags}
latches. Bit 3, TX idle, is \code{tx_empty_s and not tx_busy}, from the TX FIFO and the
transmitter. Bits 31 to 4 are zero.

None of them can be a stored register the datapath writes, because each is a copy of state that
already lives somewhere else: the FIFOs' counts, the error latches, the transmitter's state machine.
A stored copy would have to be updated by every event that changes its source, on the same edge,
by logic that duplicates the FIFO's own; any cycle in which it lagged would be a cycle in which the
driver could read ``TX ready'' for a FIFO that had just filled and push a byte into it that is then
dropped. Computed combinationally, each bit is correct by construction, and the bridge's latch-once
read takes a coherent snapshot of all four at once.

\ans{d} It relies on the bridge raising \code{reg_write} only for a \textbf{completed, non-aborted
write}: one clock, once, when the fifth byte arrives with \code{SS} still low, with \code{reg_addr}
and \code{reg_wdata} already holding the transaction's index and all four data bytes. The provided
\code{spi_reg_bridge} does exactly that: it counts the bytes, resets the count whenever \code{SS}
rises, and asserts \code{reg_write} only on byte four of a write command, which
\code{spi_reg_bridge_tb} checks, abort included.

If a half-finished transaction could reach the bank as a \code{reg_write}, its data word would be
partly stale, the bytes left over from an earlier transaction, and the bank would act on it: a
\code{BAUD_DIV} of garbage, and the line at a rate nobody chose; a \code{TX_DATA} push of a byte the
master never finished sending, which then goes out on \code{tx}; an \code{RX_POP} that discards an
unread byte; an \code{ERROR_FLAGS} write that clears an error nobody saw. Every one of those is a
single-cycle action with nothing to undo it, which is why the guarantee has to hold before the
strobe reaches the bank.

\solution{c5:ex:2}{\code{uart_top}}
\ans{b} A release that lands near a clock edge would leave some flip-flops out of reset on that edge
and others on the next, with any of them possibly metastable, and the peripheral would start
half-reset. Here that could mean a FIFO whose \code{head} leaves reset a cycle before its
\code{count}, a receiver whose \code{rx_prev} leaves reset after its state and sees a falling edge
that never happened, or the transmitter's state machine out of step with its \code{tx} register:
states the logic was never designed to leave. Asserting asynchronously is nonetheless right, because
entering reset has no such hazard. Every flop goes to its reset value immediately, together, and
without needing a clock, so the design is safe the instant \code{reset_n} falls, even if the clock
has stopped, and \code{tx} goes to its idle high at once. Only the release has to line up with the
clock, and \code{reset_sync} lines it up.

\ans{c} The port that settles it is \code{sync}'s \code{reset_s2_n} input. To make
\code{reset_s2_n} with a \code{sync} instance, you would have to connect the instance's own output,
\code{reset_s2_n}, to its own reset input: the synchronizer would be reset by the signal it has not
produced yet, which is undefined until it has been clocked, and it would never assert while it was
holding itself in reset. The raw \code{reset_n} would have to go in through \code{async_in}, so even
the assertion would wait two clock edges and would never come at all without a clock. The circle
cannot be closed; \code{reset_sync} breaks it by taking the raw reset on its flops' asynchronous
clear.

\ans{d} \code{tx_load} is \code{(not tx_empty) and (not tx_busy)}, and \code{tx_pop} is the same
signal.
\begin{itemize}
  \item \textbf{Never zero.} Whenever a byte is queued and the transmitter is idle, \code{tx_load} is
    already high, combinationally, so the transmitter loads on the next edge.
  \item \textbf{Exactly one.} On that edge, \code{uart_tx} sees \code{start} in \code{STATE_IDLE},
    latches the front byte and moves to \code{STATE_SEND}, and the FIFO pops the same byte. Because
    \code{busy} is decoded combinationally from the state register, it is high from that edge on,
    so \code{tx_load} falls in the very next cycle. It cannot reload while the transmitter is busy,
    and the FIFO cannot pop a second byte for the same frame.
  \item \textbf{Held low} for the rest of the frame by \code{tx_busy}, which stays high from the
    load edge until the edge on which the stop bit's sixteenth tick returns the state machine to
    \code{STATE_IDLE}. In that first idle cycle, a queued byte raises \code{tx_load} again.
\end{itemize}
Simulated with three bytes queued at once, the feeder popped exactly three times, once per frame,
each pop on the first idle edge after the previous frame's \code{done}. The argument depends on
\code{busy} rising on the load edge itself. With \code{busy} registered one cycle late, as a second
flip-flop, \code{uart_tx_tb} and \code{uart_top_tb} both still pass, since each sends one byte and a
second pop from an empty FIFO is ignored; but with three bytes queued the feeder popped twice on
consecutive edges, the transmitter took the first and ignored the second, and the third byte never
arrived.

\ans{e} One byte, around the loop:
\begin{enumerate}
  \item The testbench plays a five-byte write of \code{TX_DATA} on \code{sclk}, \code{mosi} and
    \code{ss}. \code{spi_slave} synchronizes the three lines and shifts each byte in;
    \code{spi_reg_bridge} assembles the command and the four data bytes and, on the fifth, raises
    \code{reg_write} with \code{reg_addr} 3 and the byte in \code{reg_wdata(7 downto 0)}.
  \item \code{uart_regs} decodes that as \code{tx_push}, and its TX \code{fifo} stores the byte.
  \item The \textbf{feeder} in \code{uart_top} sees the FIFO not empty and the transmitter not busy,
    and raises \code{tx_load}: \code{uart_tx} latches the byte into its frame, and the TX FIFO pops.
  \item \code{uart_tx} drives the frame out on \code{tx}, one bit per sixteen \code{baud_gen}
    ticks.
  \item The testbench's loopback, \code{rx <= tx}, returns it on the \code{rx} pin.
  \item \code{sync} carries \code{rx} into the clock domain as \code{rx_s2}.
  \item \code{uart_rx}, on the same ticks, finds the start edge, samples the eight bits and the stop
    bit, and pulses \code{valid}, which is \code{rx_push}.
  \item \code{uart_regs} pushes the byte into its RX \code{fifo}, and \code{STATUS} bit 1 rises.
  \item The testbench polls \code{STATUS} over SPI until it sees bit 1, then reads \code{RX_DATA}:
    \code{spi_reg_bridge} latches \code{reg_rdata}, which the bank's read mux fills from the RX
    FIFO's front, and \code{spi_slave} shifts it out on \code{miso}, most significant byte first.
\end{enumerate}
The system testbench stops there, with the byte still in the RX FIFO: it never writes \code{RX_POP}.
